\section{$R$-linear algebra}

A lot of linear algebra still works for $R$-modules, but not all of it. I will survey some of the main constructions and main results of the theory.

The most striking result in all of module theory is that if $R$ is a random ring, then there is no requirement for an $R$-module $M$ to have a basis! We can already see a simple example of this by looking at $\mathbb{Z} $-modules, which we recall (by \cref{ex:module-1}) are simply abelian groups. First, let's recall what a basis is.

\begin{definition}\label{dfn:basis}
	Let $R$ be a ring, and let $M$ be an $R$-module. A finite subset $\left\{ m_1, \ldots ,m_n \right\} \subseteq M$ is \emph{linearly independent} if the only collection of scalars $r_1, \ldots ,r_n\in R$ such that
	\[
		r_1m_1+ \cdots + r_nm_n = 0
	\]
	are $r_1, \ldots , r_n = 0$. The finite subset $\left\{ m_1, \ldots ,m_n \right\} \subseteq M$ is a \emph{spanning set} if for all $x\in M$, there are coefficients $r_1, \ldots , r_n\in R$ such that
	\[
		x = r_1m_1 + \cdots + r_nm_n
	.\]
	A \emph{basis} is a set $S\subseteq M$ that is both linearly independent and a spanning set.
\end{definition}

Consider the following very simple abelian group. Let $\mathbb{Z} /2$ be the set $\left\{ 0,1 \right\} $ with addition defined by
\begin{align*}
	&0+0=0&0+1=1 \\
	&1+0=1&1+1=0
\end{align*}
The only difference from normal addition is $1+1=0$. We call this the group of \emph{integers modulo 2}. Other notations for this group are $\mathbb{Z} /(2)$, $\mathbb{Z} /2\mathbb{Z} $, and $\mathbb{Z}_2$. The addition operation is typically written as $n+m\pmod{2}$, but we will drop the $(\text{mod}\;2)$.

Since all abelian groups are $\mathbb{Z} $-modules, we see that $\mathbb{Z} /2$ is a $\mathbb{Z} $-module. Does this $\mathbb{Z} $-module have a basis? There are only 4 subsets of $\mathbb{Z} /2$: the empty set, and the sets $\left\{ 0\right\} ,\left\{ 1 \right\} ,\left\{ 0,1 \right\} $. Clearly the empty set is not spanning. Additionally, $\left\{ 0 \right\} $ is not spanning, since there is no way to add 0 to itself in order to get 1. The set $\left\{ 1 \right\}$ \emph{is} spanning, since $2\cdot 1 = 1+1 = 0$. However, it's \emph{not linearly independent}. By the equation we just wrote, $2\cdot 1=0$ for $2\in\mathbb{Z} $. However, since $2\neq 0$ in $\mathbb{Z} $, this immediately fails linear independence. The set $\left\{ 0,1 \right\} $ is also not linearly independent for the same reason: $0\cdot 0+2\cdot 1=0$, but $2\neq 0$.

This immediately makes linear algebra with $R$-modules much harder. For example, without a basis, we can't define the \emph{dimension} of an $R$-module. In order to talk about matrices, we need to have a basis; thus we cannot always do matrix algebra in arbitrary $R$-modules.

\begin{definition}\label{dfn:free-module}
	An $R$-module $M$ with a basis $S\subseteq M$ is called a \emph{free $R$-module}. The basis is also referred to as the \emph{generating set} of $M$.
\end{definition}

Although not all $R$-modules are free (unlike vector spaces), we can still do quite a lot of linear algebra with $R$-modules. Let's introduce some notions from linear algebra that extend to $R$-modules without much difficulty.

\begin{definition}\label{dfn:R-linear}
	Let $R$ be a ring, and let $M,N$ be $R$-modules. An \emph{$R$-linear map} from $M$ to $N$ is a function $\varphi : M \to N$ with the following properties:
	\begin{enumerate}
		\item For all $m_1,m_2\in M$, we have $\varphi (m_1+m_2) = \varphi (m_1)+\varphi (m_2)$.
		\item For all $m\in M$ and $r\in R$, we have $\varphi (r m)=r\cdot \varphi (m)$.
	\end{enumerate}
\end{definition}
Note that \cref{dfn:R-linear} is the same definition as a normal linear map, except the scalars now come from a ring $R$, rather than $\mathbb{R} $ or $\mathbb{C} $.

\begin{example}\label{ex:R-linear-1}
	Let $M,N$ be $R$-modules. The map $0: M \to N$ which is defined by $0(m)=0$ is an $R$-linear map. We call this the \emph{zero morphism}, the \emph{zero map}, or the \emph{trivial map}. Another example is the map $1_M : M \to M$ defined by $1_M(m) = m$. This is called the \emph{identity morphism} of $M$.
\end{example}

\begin{example}\label{ex:R-linear-2}
	Every linear map between vector spaces you have ever seen is an $\mathbb{R} $-linear (or $\mathbb{C} $-linear) map.
\end{example}

We can similarly extend a lot of definitions like vector subspaces, kernels and images to the setting of $R$-modules.

\begin{definition}\label{dfn:submod-kernel-image}
	Let $R$ be a ring, and let $M,N$ be $R$-modules.
	\begin{enumerate}
		\item We say $N$ is a \emph{submodule} of $M$ if $N \subseteq M$ as a set, and if $N$ is closed under addition and scalar multiplication.
		\item Let $\varphi : M \to N$ be an $R$-linear map. The \emph{kernel} of $\varphi $, denoted $\Ker \varphi $, is the submodule
		\[
			\Ker \varphi =\left\{ m\in M\mid \varphi (m)=0 \right\} \subseteq M
		.\]
	\item Let $\varphi : M \to N$ be an $R$-linear map. The \emph{image} of $\varphi $, denoted $\Im \varphi $, is the submodule
	\[
		\Im \varphi =\left\{ \varphi (m)\in N\mid m\in M \right\} \subseteq N
	.\]
	\end{enumerate}
\end{definition}

We can also define Cartesian products (which we simply refer to as products), which will be extremely useful.

\begin{definition}\label{dfn:product}
	Let $R$ be a ring, and let $M$ and $N$ be $R$-modules. We define the \emph{product module} $M\times N$ to be the $R$-module defined as follows.
	\begin{itemize}
		\item The underlying set $M\times N$ consists of all ordered pairs $(m,n)$, where $m\in M$ and $n\in N$.
		\item The addition operation is defined \emph{pointwise}, meaning for all $m_1,m_2\in M$ and $n_1,n_2\in N$, we have
		\[
			(m_1,n_1) + (m_2,n_2) = (m_1 + m_2,n_1+n_2)
		.\]	
		\item The multiplication operation is also defined pointwise, meaning for all $m\in M$, $n\in N$, and $r\in R$, we define
		\[
			r(m,n) = (r m,rn)
		.\]
	\end{itemize}
\end{definition}

\begin{notation}\label{notation:direct-sum}
	Due to some compliated but very nice properties of $R$-modules, the notation $M\times N$ for the product is not used very often. Most of the times, we write $M\times N$ as $M\oplus N$, and call it either the \emph{product} or the \emph{direct sum}. The element $(m,n)$ is often written as $m+n$. This notation makes sense since we can write
	\[
		(m,n) = (m,0) + (0,n),
	\]
	and then we simply use the notation $m=(m,0)$ and $n=(0,n)$. This notation turns out to be extremely convenient, and is used extremely often. To avoid confusion, we will not be using this notation (we will be using ordered pairs and $\times $), because this is more aligned with undergraduate linear algebra notation. However, it's important to know that this notation exists.
\end{notation}

\begin{remark}\label{rmk:prod-coprod}
	The product of $R$-modules is both a product \emph{and} a coproduct, in the sense of category theory.
\end{remark}

We will now discuss something known as \emph{torsion}, which can make a module fail to have a basis. Let $M$ be an $R$-module, and let $r\in R$ and $m\in M$ such that $r m=0$. Suppose that $r\neq 0$. Does $m=0$? If $R=\mathbb{R} $, then the answer is yes, since we can simply multiply by $1/r$ (since $r\neq 0$) to get
\[
	0=\frac{1}{r} r m = \frac{r}{r} m = 1m = m
.\]
However, if $R$ is an arbitrary ring, we can't do this. For example, consider the ring of $n\times n$ matrices with real entries $M_n(\mathbb{R} )$. Recall that \emph{not all matrices are invertible}. Thus, if $r\in M_n(\mathbb{R} )$ is some matrix, there might not be some element $r^{-1} $ allowing us to cancel $r$ in the equation $r m=0$.

\begin{definition}\label{dfn:torsion}
	Let $R$ be a ring, and $M$ an $R$-module. An element $m\in M$ is called \emph{torsion} if there is some nonzero $r\in R$ such that $r m=0$. The \emph{torsion submodule} of $M$, denoted $\operatorname{Tor} (M)$, is the set
	\[
		\operatorname{Tor} (M) = \left\{ m\in M\mid m\text{ is torsion}  \right\} 
	.\]
\end{definition}

It turns out that torsion is the exact obstruction to being free. The only torsion element of a vector space is 0, so if $V$ is a vector space, then $\operatorname{Tor} (V)=\left\{ 0 \right\} $. However, there can be $R$-modules with nonzero torsion elements.

For example, consider the $\mathbb{Z} $-module $\mathbb{Z} /2$ from before, which we know has no basis. The element 0 is obviously torsion, so let's consider the element 1. Recall that $2\cdot 1=0$, but $2\neq 0$. Thus, $1$ is torsion. We thus have $\mathbb{Z} /2 = \operatorname{Tor} (\mathbb{Z} /2)=\left\{ 0,1 \right\} $.

Note that the submodule $\operatorname{Tor} (M)$ \emph{cannot} be free: given any element $m\in \operatorname{Tor} (M)$, there is some nonzero $r\in R$ with $r m=0$. Thus, any collection $m_1, \ldots ,m_n$ of torsion elements admits nonzero scalars $r_1, \ldots ,r_n\in R$ such that $r_i m_i=0$ for all $i$, and thus
\[
	r_1m_1 + \cdots + r_nm_n=0
.\]
Thus, no nonempty subset of $\operatorname{Tor} (M)$ is linearly independent.

This turns out to be an error with the ring $R$ itself: not all scalars are invertible. Thus, we can look for examples of rings where all (nonzero) elements are invertible.

\begin{definition}\label{dfn:field}
	A \emph{field} $\mathbb{K} $ is a ring $(\mathbb{K} ,+,\cdot )$ such that the following additional properties hold.
	\begin{enumerate}
		\item $\mathbb{K} $ is a commutative ring, i.e. for all $a,b\in\mathbb{K} $, we have $ab=ba$.
		\item For all nonzero $a\in\mathbb{K} $, there is some element $a^{-1} \in\mathbb{K} $ such that $a a^{-1} =a^{-1} a=1$.
	\end{enumerate}
\end{definition}

The main three examples of fields are the rational numbers $\mathbb{Q} $, the real numbers $\mathbb{R} $, and the complex numbers $\mathbb{C} $. There are other interesting examples of fields, but we will only consider these three for simplicity.

\begin{definition}\label{dfn:vector-space}
	A \emph{vector space} is a module over a field. If $\mathbb{K} $ is a field, a module over $\mathbb{K} $ is called a \emph{$\mathbb{K} $-vector space}.
\end{definition}

The final thing I want to define before we continue past $R$-modules is the notion of a \emph{quotient}. In higher mathematics, quotients are a ubiquetous construction, but they can be complex to understand at first. The idea is slightly easier to understand in the context of $R$-modules (and most other algebraic structures) than it is for most objects, like sets or topological spaces. We will begin with some motivation.

Suppose $M$ is an $R$-module, and $R$ is not a field. It would be nice if $M$ was free, but this would require that $M$ has no torsion, and we cannot know in general if $M$ has no torsion. So we will do the next best thing: given a module $M$ with some nontrivial torsion $\operatorname{Tor} (M) \neq \left\{ 0 \right\} $, how do we turn $M$ into a new module, which I will call $M / \operatorname{Tor} (M)$, which has no torsion?

Recall that a torsion element of $M$ is some $m\in M$ such that there is $r\in R$, $r\neq 0$ and $r m=0$. If this torsion element $m$ is also nonzero, then this is an obstruction to linear independence in $M$. Thus, we would like the only torsion element to be 0. In other words, if $m$ is a nonzero torsion element in $M$, we would like to somehow ``turn $m$ into 0'' when we go to the new module $M / \operatorname{Tor} (M)$. If we can turn all of the torsion elements into 0, then we have gotten rid of the obstructions to linear independence posed by the torsion elements. But how can we possibly construct such a module?

\begin{construction}\label{const:quotient}
	Let $M$ be an $R$-module, and let $N\subseteq M$ be a \emph{submodule}. In the previous motivational section, we took $N = \operatorname{Tor} (M)$, but this construction works more generally, so here we allow $N$ to be any submodule. We will define a new $R$-module $M/N$ as follows. First, let $m\in M$. We will use the notation $m+N$ to denote the following \emph{set}:
	\[
		m+N \coloneqq \left\{ m+n \mid n\in N \right\} \subseteq M
	.\]
	We call $m+N$ a \emph{coset}. We will define $M/N$ to be the collection of all cosets. To define addition, let $m_1,m_2\in M$. Then we define
	\[
		(m_1 + N) + (m_2 + N) \coloneqq (m_1+m_2)+N
	.\]
	Since $m_1+m_2\in M$, we have that $(m_1+m_2)+N$ is a coset. Additionally, since this set consists of all elements $m_1+m_2+n$, we can prove that $M/N$ is an abelian group by simply applying the fact that $M$ is an abelian group to the general form of each element of a coset. For example, to prove that addition is commutative, observe that if $m_1 + m_2 + n\in (m_1+m_2)+N$, then
	\[
		m_1 + m_2 + n = m_2 + m_1 + n \in (m_2 + m_1) + N
	.\]
	Thus, $(m_1+m_2)+N \subseteq (m_2 + m_1) + N$. The inclusion in the reverse direction is proven in the exact same way.

	To define scalar multiplication, if $r\in R$ and $m\in M$, we define
	\[
		r(m+N) \coloneqq (r m)+N
	.\]
	Since $r m\in M$, we have that $(r m)+N$ is indeed a coset, and we write it as $r m+N$. The rest of the module axioms are proven similarly to the abelian group axioms.
\end{construction}

\begin{definition}\label{dfn:quotient}
	Let $M$ be an $R$-module, and $N\subseteq M$ a submodule. The module $M/N$ defined in \cref{const:quotient} is called the \emph{quotient} of $M$ by $N$.
\end{definition}

We have now constructed a new module, but that isn't quite good enough. We wanted to \emph{turn $M$ into a new module}. We have that new module now, but we don't have a process for taking an element of $m$, and turning it into an element of $M/N$. For this, we need the \emph{natural projection}.

Given any submodule $N\subseteq M$, there is an $R$-linear map $\pi : M \to M/N$ called the \emph{natural} or \emph{canonical projection} (also simply referred to as the quotient, or the quotient morphism). This map is defined as follows: if $m\in M$, then
\[
	\pi (m) = m+N
.\]
We can prove that this is $R$-linear. First, if $m_1,m_2\in M$, then
\[
	\pi (m_1) + \pi (m_2) = m_1 + N + m_2 + N = (m_1 + m_2) + N = \pi (m_1 + m_2)
.\]
Next, if $m\in M$ and $r\in R$, then
\[
	r\pi (m) = r(m+N) = (r m)+N = \pi (r m)
.\]
This is the tool we needed for taking elements from $M$ to elements of $M/N$.

Now let's return to our investigation of $M / \operatorname{Tor} (M)$. Given an element $m\in \operatorname{Tor} (M)$, we want this element to be sent to 0 in the quotient module. The zero element of $M / \operatorname{Tor} (M)$ is the coset $0+\operatorname{Tor} (M)$. If $m$ is torsion, then we simply need to show that $\pi (m) = m+\operatorname{Tor} (M)$ is equal to $0+\operatorname{Tor} (M)$.

First, note that $m\in \operatorname{Tor} (M)$. Since $\operatorname{Tor} (M)$ is a submodule of $M$, it is in particular a module, and thus contains $-m\in \operatorname{Tor} (M)$. Therefore, the element $m+(-m) = 0$ is in the coset $m+\operatorname{Tor} (M)$. Thus, if I have some element $0+t\in 0+\operatorname{Tor} (M)$, then I can rewrite
\[
	0+t = (m+(-m))+t = m+(-m + t)\in m + \operatorname{Tor} (M)
.\]
Here, we note that $-m+t\in \operatorname{Tor} (M)$ since $-m,t\in \operatorname{Tor} (M)$, and since $\operatorname{Tor} (M)$ is a module, it is closed under addition. Thus, $0+\operatorname{Tor} (M) \subseteq m+ \operatorname{Tor} (M)$.

The other inclusion is similar. Let $m+t\in \operatorname{Tor} (M)$. Then since $-m,t\in \operatorname{Tor} (M)$, we have $-m+t\in \operatorname{Tor} (M)$, and thus $m+(-m+t) = 0+t\in m+\operatorname{Tor} (M)$. Since $0+t\in 0+\operatorname{Tor} (M)$, we conclude that $m+\operatorname{Tor} (M)\subseteq 0+\operatorname{Tor} (M)$. Thus,
\[
	m+\operatorname{Tor} (M) = 0+\operatorname{Tor} (M)
.\]

Finally, note that since $0\in \operatorname{Tor} (M)$, if an element $\widetilde{m} \in M$ is \emph{not} torsion, then $\widetilde{m} +0 = \widetilde{m} \in \widetilde{m} + \operatorname{Tor} (M)$. Since $\widetilde{m} \notin \operatorname{Tor} (M)$, we can conclude that $\widetilde{m} +\operatorname{Tor} (M) \neq 0+\operatorname{Tor} (M)$, and thus $\widetilde{m} $ is nonzero. Therefore, the \emph{only} elements of $M$ that are sent to 0 in $M/\operatorname{Tor} (M)$ are the torsion elements. This produces a new module without torsion. It can in fact be shown that this module is free, although that would take more time than we have.

The idea of quotients is very useful. Although they are very abstract constructions, oftentimes they can be easier to understand than some very complex modules (and some special quotients have nicer descriptions than cosets). Thus, being able to reduce a very complex module to a quotient of a more simple module can help us understand it more. I will introduce two extremely powerful theorems of quotients (without proof) to try and illustrate this idea.

\begin{theorem}[First isomorphism theorem]\label{thm:first-iso}
	Let $R$ be a ring, $M,N$ modules, and $\varphi : M \to N$ some $R$-linear map. Then there is an isomorphism $\widetilde{\varphi } : M / \Ker \varphi  \to \Im \varphi $, defined by
	\[
		\widetilde{\varphi } (m+\Ker \varphi ) \coloneqq \varphi (m)
	.\]
	In particular, if $\varphi $ is surjective, then $M / \Ker \varphi \cong N$.
\end{theorem}

\begin{theorem}\label{cor:dimension}
	Let $V$ be a finite dimensional vector space over some field $\mathbb{K} $, and let $W\subseteq V$ be a subspace. Then
	\[
		\dim (V/W) = \dim (V) - \dim (W)
	.\]
\end{theorem}
